\documentclass[10pt]{article} 
\usepackage{listings}
\usepackage{amssymb}
\usepackage{alltt}
\begin{document}
\lstset{language=Fortran}
\setlength\parindent{0pt}
%*************************
%*************************
%*************************
\section{Introduction}

%*************************
\subsection{What is KROME?}
KROME is a nice and friendly chemistry package for a wide range of astrophysical simulations. Given a chemical
network (in CSV format) it automatically generates all the routines needed to solve the kinetic of the system, 
modeled as system of coupled Ordinary Differential Equations. It provides different options which make it unique and very
flexible. The following user-guide is aimed at providing an easy approach to the package. Any suggestions and
comments are welcomed. KROME is an open-source package, GNU-licensed, and any improvements provided by the
users is well accepted.

%*************************
\subsection{Why KROME?}
Chemistry plays a key role\dots

%*************************
%*************************
%*************************
\section{Getting started}

%*************************
\subsection{Install KROME}
To run KROME you need only to download the package and unpack in your favorite folder.

%*************************
\subsection{Prerequisites}
KROME has been developed to work on linux systems, but can run also on different OS if prerequisites are satisfied.
KROME needs a fortran compiler as GNU gfortran or Intel IFORT (suggested), and python (>=2.5).
Optional: if you want to graphically display the results from the default tests provided with the package, 
we suggest to install gnuplot (>=3.7.1). Each test has its gnuplot script (.gps) aimed at showing the output properly.

%*************************
\subsection{Quickstart}
%--------------------------
\subsubsection{Overview}
In this section we show the main features of KROME needed for a quickstart.
The scheme is quite simple: 
\begin{equation}
	\mathrm{NETWORK\,\,FILE}\longrightarrow\mathrm{KROME}\longrightarrow\mathrm{F90\,\,ROUTINES}\,.\nonumber
\end{equation}
First, you need a NETWORK FILE that describes the network (see \ref{sect:input_file_specifications} for details). With this input KROME provides all the F90 ROUTINES needed for computing the chemical evolution in your code.

The interface between KROME and your code is simply given by a call to
\begin{alltt}
	call krome(x(:), temperature, density, time_step)
\end{alltt}
where the arguments are $x(:)$ the array containing the initial abundances of the species, the $temperature$ and $density$ of the gas, and finally the $time step$ of the total integration. The output of final abundances is the array $x(:)$.

%--------------------------
\subsubsection{A simple example}
Here an example of KROME.

We use a simple kinetic system described by three reactions
\begin{equation}
\mathrm{CH}+\mathrm{H} \leftrightharpoons \mathrm{CH_2} \rightarrow \mathrm{H_2}+\mathrm{C}\,.
\end{equation}
They are controlled by three rate coefficients $k_1=4\times10^{-10}\,T_g$, $k_2=10^{-10}\,T_g$, $k_1=3\times10^{-9}\,T_g$,
and each reaction is valid into the temperature range $T_g=[10,100]$ K. These information are represented in the network file as
\begin{alltt}
	1,CH,H,,CH2,,,,10.,100.,4.d-10*Tgas
	2,CH2,,,CH,H,,,10.,100.,1.d-10*Tgas
	3,CH2,,,C,H2,,,10.,100.,3.d-9*Tgas
\end{alltt}

and this file is located in the folder \verb|networks/react_simple|.

At the beginning of our calculations all the species are set to zero except CH and H that are $x_\mathrm{CH}=x_\mathrm{H}=0.5$ expressed as mass fractions (i.e. it is a gas half in mass hydrogen and half CH).

We want to determine the evolution of all the species in gas during time, and to do that one have to:
\begin{enumerate}
	\item Go into KROME directory and type
	\begin{alltt}
		> python krome.py networks/react_simple
	\end{alltt}
	this will create a set of f90 files in the folder \verb|build/| that are the files containing all the subroutines to compute the evolution. At this stage you don't need to check inside these files.

	\item Now copy the Makefile to the build directory using
	\begin{alltt}
		> cp tests/Makefile build/.
	\end{alltt}
	in this way you have an example Makefile that can be modified to your purposes.

	\item The main f90 program is obtained by using
	\begin{alltt}
		> cp tests/test.f90 build/.
	\end{alltt}
	This file is a simple f90 program that call KROME to calculate the evolution of the species. In this example test.f90 acts as a global code that needs a subroutine to calculate the chemical evolution. This file is widely discussed in section 
\ref{sssect:detail_simple_example}, please have a look at this.

	\item Compile the code by 
	\begin{alltt}
		> cd build
		> make
	\end{alltt}
	It will provide a \verb|krome| executable.

	\item Run the test using 
	\begin{alltt}
		> ./krome > out
	\end{alltt}
	redirecting the output to the file \verb|out|.
	
	\item Now you can plot the results using a gnuplot script we provide. First copy it into the current folder
	\begin{alltt}
		> cp ../tests/plots/plot.gps .
	\end{alltt}
	Then using gnuplot you obtain the evolutionary tracks of the different species
	\begin{alltt}
		> gnuplot
		gnuplot> load 'plot.gps'
	\end{alltt}
	In this example you see that CH is destroyed and the intermediate species CH$_2$ (blue) first increases its content due to the reaction $k_1$, and than it decreases because of $k_3$ that has the largest rate coefficient. At the end of the evolution
C and H$_2$ are formed. The reason why H abundance is still high is that H has no more CH reaction partners to react with, in fact setting
$x_\mathrm{CH}=x_\mathrm{H}=0.5$ will produce a number density ratio of $n_\mathrm{H}/n_\mathrm{CH}=7$ due to their mass difference.
\end{enumerate}

%--------------------------
\subsubsection{Simple example in detail}\label{sssect:detail_simple_example}
Details of test.f90

%*************************
\subsection{Input File Specifications}\label{sect:input_file_specifications}
Input File Specifications

%*************************
%*************************
%*************************
\section{Tests}
%*************************
\subsection{TEST1: Bare-minimal test}
A tiny little test.
%*************************
\subsection{TESTDark Cloud test}
WH2008 test.


%*************************
%*************************
%*************************
\section{Theoretical Background}


\end{document} 
